%!TEX root = ../template.tex
%%%%%%%%%%%%%%%%%%%%%%%%%%%%%%%%%%%%%%%%%%%%%%%%%%%%%%%%%%%%%%%%%%%%
%% abstract-en.tex
%% ISEL thesis document file
%%
%% Abstract in English
%% (Replace this text with your own abstract)
%%%%%%%%%%%%%%%%%%%%%%%%%%%%%%%%%%%%%%%%%%%%%%%%%%%%%%%%%%%%%%%%%%%%
\textbf{This document is both the ISEL \LaTeX{} template for bachelor's project reports and master's dissertations, and its user manual. If you only read this page, read the three rules below:}

\bigskip
\noindent%
\fcolorbox{iselorange}{white}{%
\begin{minipage}{\dimexpr\linewidth-2\fboxsep-2\fboxrule}
\begin{enumerate}
	\item \textbf{Compile with LuaLaTeX}, not pdfLaTeX. Only then is the Verdana font required by the ISEL guidelines used.
	\item \textbf{Edit only} the \texttt{Config} folder (options, cover and list of files), the \texttt{Chapters} folder (text and images) and the \texttt{bibliography.bib} file.
	\item \textbf{Do not change} \texttt{template.tex}, \texttt{iselthesis.cls} or the \texttt{ISELthesis-files} folder.
\end{enumerate}
\end{minipage}%
}

\bigskip
Details are in Chapter~\ref{cha:users_manual} (structure and options) and Chapter~\ref{cha:a_short_latex_tutorial_with_examples} (compiling and examples).

\textbf{About the abstract.} The template places the abstract in the main language of the document first, followed by the abstract in the other language (English or Portuguese). The abstract must not exceed \textbf{one page or 300 words}, and should answer four questions: \textbf{what is the problem}, \textbf{why is it interesting}, \textbf{what is the solution}, and \textbf{what follows from the solution}.

% Keywords of the abstract in English
\begin{keywords}
Keywords (in English) \ldots -- at least 3 (three)
\end{keywords}